\documentclass[11pt]{article}
\usepackage[review]{acl}

\usepackage{times}
\usepackage{latexsym}
\usepackage[T1]{fontenc}
\usepackage[utf8]{inputenc}
\usepackage{microtype}
\usepackage{inconsolata}
\usepackage{graphicx}
\usepackage{booktabs}
\usepackage{array}
\usepackage{xcolor}
\usepackage{amsmath}
\usepackage{amssymb}
\usepackage{tikz}
\usetikzlibrary{arrows.meta,shapes.geometric}
\usepackage{url}

% Auto-generated by figures/make_figures.py
\newcommand{\NQueries}{100}
\newcommand{\NProviders}{3}
\newcommand{\NTraces}{300}
\newcommand{\NJudgeTotal}{6,909}
\newcommand{\NJudgeValid}{6,869}
\newcommand{\NJudgeInvalid}{40}
\newcommand{\NJudgeSnippetValid}{6,519}
\newcommand{\NJudgePageValid}{350}
\newcommand{\AllCorrect}{10}
\newcommand{\TwoCorrect}{12}
\newcommand{\OneCorrect}{22}
\newcommand{\AllWrong}{56}
\newcommand{\OneOrMoreCorrect}{44}
\newcommand{\MetricMode}{raw_judge+semantic_tsv}
\newcommand{\CorrectMetricName}{semantic-correct}
\newcommand{\BraveEM}{21}
\newcommand{\BraveCorrect}{25}
\newcommand{\BraveCorrectGain}{4}
\newcommand{\BraveFone}{0.270}
\newcommand{\BraveSearchAvg}{2.29}
\newcommand{\BraveFetchAvg}{1.02}
\newcommand{\BraveFetchedPct}{65\%}
\newcommand{\BraveAvgTokens}{59,627}
\newcommand{\BraveTokensM}{5.96}
\newcommand{\BraveMedianTokens}{40,638}
\newcommand{\BraveMaxTokens}{303,462}
\newcommand{\BraveOverHundredK}{19}
\newcommand{\BraveFetchSuccess}{92}
\newcommand{\BraveFetchFailed}{10}
\newcommand{\BraveAnswered}{98}
\newcommand{\BraveAbstained}{2}
\newcommand{\BraveGoldURLExact}{59}
\newcommand{\BraveGoldDomain}{82}
\newcommand{\BraveGoldFamily}{61}
\newcommand{\BraveSnippetSurfaceHit}{71}
\newcommand{\BraveAnswerPage}{52}
\newcommand{\BraveAnswerAvailable}{78}
\newcommand{\BraveWrongWithAnswer}{57}
\newcommand{\BraveWrongWithoutAnswer}{22}
\newcommand{\BravePreFetchSupportQ}{30}
\newcommand{\BravePostFetchDiscoveredSupportQ}{3}
\newcommand{\BraveTrajectoryVisibleSupportQ}{33}
\newcommand{\BraveNoPreFetchSupportQ}{70}
\newcommand{\BraveSupportVisibleQ}{33}
\newcommand{\BraveNoSupportQ}{70}
\newcommand{\BraveSnippetRows}{2,095}
\newcommand{\BravePageRows}{101}
\newcommand{\BraveGoldRows}{97}
\newcommand{\BravePreFetchSupportRows}{101}
\newcommand{\BravePostFetchDiscoveredRows}{3}
\newcommand{\BraveContraRows}{89}
\newcommand{\BraveContraRatio}{0.92}
\newcommand{\BraveRankOneGold}{13}
\newcommand{\BraveRankOnePct}{13\%}
\newcommand{\BraveSmartN}{3}
\newcommand{\BraveSmartCorrect}{3}
\newcommand{\BraveSmartEM}{3}
\newcommand{\BraveSmartRate}{100\%}
\newcommand{\BraveSmartCI}{[44--100]}
\newcommand{\BraveMissedN}{27}
\newcommand{\BraveMissedCorrect}{11}
\newcommand{\BraveMissedEM}{11}
\newcommand{\BraveMissedRate}{41\%}
\newcommand{\BraveMissedCI}{[25--59]}
\newcommand{\BraveBlindN}{54}
\newcommand{\BraveBlindCorrect}{11}
\newcommand{\BraveBlindEM}{11}
\newcommand{\BraveBlindRate}{20\%}
\newcommand{\BraveBlindCI}{[12--33]}
\newcommand{\BraveNoopN}{16}
\newcommand{\BraveNoopCorrect}{0}
\newcommand{\BraveNoopEM}{0}
\newcommand{\BraveNoopRate}{0\%}
\newcommand{\BraveNoopCI}{[0--19]}
\newcommand{\BraveSmartMinusMissed}{+59}
\newcommand{\TavilyEM}{21}
\newcommand{\TavilyCorrect}{25}
\newcommand{\TavilyCorrectGain}{4}
\newcommand{\TavilyFone}{0.261}
\newcommand{\TavilySearchAvg}{2.74}
\newcommand{\TavilyFetchAvg}{1.30}
\newcommand{\TavilyFetchedPct}{76\%}
\newcommand{\TavilyAvgTokens}{54,156}
\newcommand{\TavilyTokensM}{5.42}
\newcommand{\TavilyMedianTokens}{36,867}
\newcommand{\TavilyMaxTokens}{305,474}
\newcommand{\TavilyOverHundredK}{16}
\newcommand{\TavilyFetchSuccess}{119}
\newcommand{\TavilyFetchFailed}{11}
\newcommand{\TavilyAnswered}{97}
\newcommand{\TavilyAbstained}{3}
\newcommand{\TavilyGoldURLExact}{57}
\newcommand{\TavilyGoldDomain}{82}
\newcommand{\TavilyGoldFamily}{60}
\newcommand{\TavilySnippetSurfaceHit}{60}
\newcommand{\TavilyAnswerPage}{57}
\newcommand{\TavilyAnswerAvailable}{75}
\newcommand{\TavilyWrongWithAnswer}{56}
\newcommand{\TavilyWrongWithoutAnswer}{23}
\newcommand{\TavilyPreFetchSupportQ}{16}
\newcommand{\TavilyPostFetchDiscoveredSupportQ}{8}
\newcommand{\TavilyTrajectoryVisibleSupportQ}{24}
\newcommand{\TavilyNoPreFetchSupportQ}{84}
\newcommand{\TavilySupportVisibleQ}{24}
\newcommand{\TavilyNoSupportQ}{84}
\newcommand{\TavilySnippetRows}{2,339}
\newcommand{\TavilyPageRows}{125}
\newcommand{\TavilyGoldRows}{31}
\newcommand{\TavilyPreFetchSupportRows}{34}
\newcommand{\TavilyPostFetchDiscoveredRows}{9}
\newcommand{\TavilyContraRows}{58}
\newcommand{\TavilyContraRatio}{1.87}
\newcommand{\TavilyRankOneGold}{17}
\newcommand{\TavilyRankOnePct}{50\%}
\newcommand{\TavilySmartN}{3}
\newcommand{\TavilySmartCorrect}{1}
\newcommand{\TavilySmartEM}{1}
\newcommand{\TavilySmartRate}{33\%}
\newcommand{\TavilySmartCI}{[6--79]}
\newcommand{\TavilyMissedN}{13}
\newcommand{\TavilyMissedCorrect}{7}
\newcommand{\TavilyMissedEM}{7}
\newcommand{\TavilyMissedRate}{54\%}
\newcommand{\TavilyMissedCI}{[29--77]}
\newcommand{\TavilyBlindN}{63}
\newcommand{\TavilyBlindCorrect}{16}
\newcommand{\TavilyBlindEM}{16}
\newcommand{\TavilyBlindRate}{25\%}
\newcommand{\TavilyBlindCI}{[16--37]}
\newcommand{\TavilyNoopN}{21}
\newcommand{\TavilyNoopCorrect}{1}
\newcommand{\TavilyNoopEM}{1}
\newcommand{\TavilyNoopRate}{5\%}
\newcommand{\TavilyNoopCI}{[1--23]}
\newcommand{\TavilySmartMinusMissed}{-21}
\newcommand{\FirecrawlEM}{23}
\newcommand{\FirecrawlCorrect}{26}
\newcommand{\FirecrawlCorrectGain}{3}
\newcommand{\FirecrawlFone}{0.282}
\newcommand{\FirecrawlSearchAvg}{2.51}
\newcommand{\FirecrawlFetchAvg}{1.28}
\newcommand{\FirecrawlFetchedPct}{81\%}
\newcommand{\FirecrawlAvgTokens}{57,979}
\newcommand{\FirecrawlTokensM}{5.80}
\newcommand{\FirecrawlMedianTokens}{34,383}
\newcommand{\FirecrawlMaxTokens}{380,646}
\newcommand{\FirecrawlOverHundredK}{16}
\newcommand{\FirecrawlFetchSuccess}{121}
\newcommand{\FirecrawlFetchFailed}{7}
\newcommand{\FirecrawlAnswered}{96}
\newcommand{\FirecrawlAbstained}{4}
\newcommand{\FirecrawlGoldURLExact}{60}
\newcommand{\FirecrawlGoldDomain}{82}
\newcommand{\FirecrawlGoldFamily}{64}
\newcommand{\FirecrawlSnippetSurfaceHit}{54}
\newcommand{\FirecrawlAnswerPage}{61}
\newcommand{\FirecrawlAnswerAvailable}{76}
\newcommand{\FirecrawlWrongWithAnswer}{53}
\newcommand{\FirecrawlWrongWithoutAnswer}{24}
\newcommand{\FirecrawlPreFetchSupportQ}{16}
\newcommand{\FirecrawlPostFetchDiscoveredSupportQ}{3}
\newcommand{\FirecrawlTrajectoryVisibleSupportQ}{19}
\newcommand{\FirecrawlNoPreFetchSupportQ}{84}
\newcommand{\FirecrawlSupportVisibleQ}{19}
\newcommand{\FirecrawlNoSupportQ}{84}
\newcommand{\FirecrawlSnippetRows}{2,085}
\newcommand{\FirecrawlPageRows}{124}
\newcommand{\FirecrawlGoldRows}{27}
\newcommand{\FirecrawlPreFetchSupportRows}{30}
\newcommand{\FirecrawlPostFetchDiscoveredRows}{3}
\newcommand{\FirecrawlContraRows}{70}
\newcommand{\FirecrawlContraRatio}{2.59}
\newcommand{\FirecrawlRankOneGold}{4}
\newcommand{\FirecrawlRankOnePct}{13\%}
\newcommand{\FirecrawlSmartN}{3}
\newcommand{\FirecrawlSmartCorrect}{1}
\newcommand{\FirecrawlSmartEM}{1}
\newcommand{\FirecrawlSmartRate}{33\%}
\newcommand{\FirecrawlSmartCI}{[6--79]}
\newcommand{\FirecrawlMissedN}{13}
\newcommand{\FirecrawlMissedCorrect}{7}
\newcommand{\FirecrawlMissedEM}{7}
\newcommand{\FirecrawlMissedRate}{54\%}
\newcommand{\FirecrawlMissedCI}{[29--77]}
\newcommand{\FirecrawlBlindN}{70}
\newcommand{\FirecrawlBlindCorrect}{18}
\newcommand{\FirecrawlBlindEM}{18}
\newcommand{\FirecrawlBlindRate}{26\%}
\newcommand{\FirecrawlBlindCI}{[17--37]}
\newcommand{\FirecrawlNoopN}{14}
\newcommand{\FirecrawlNoopCorrect}{0}
\newcommand{\FirecrawlNoopEM}{0}
\newcommand{\FirecrawlNoopRate}{0\%}
\newcommand{\FirecrawlNoopCI}{[0--22]}
\newcommand{\FirecrawlSmartMinusMissed}{-21}
\newcommand{\BraveSemExOneGold}{Astra Zeneca}
\newcommand{\BraveSemExOneAns}{AstraZeneca}
\newcommand{\BraveSemExOneNote}{Spacing only: "Astra Zeneca" == "AstraZeneca"}
\newcommand{\BraveSemExTwoGold}{UnionPay}
\newcommand{\BraveSemExTwoAns}{China UnionPay}
\newcommand{\BraveSemExTwoNote}{Same entity: "China UnionPay" is the full name of "UnionPay"}
\newcommand{\BraveSemExThreeGold}{US\$120,000}
\newcommand{\BraveSemExThreeAns}{\$120,000}
\newcommand{\BraveSemExThreeNote}{Same amount: "\$120,000" == "US\$120,000"}
\newcommand{\TavilySemExOneGold}{Astra Zeneca}
\newcommand{\TavilySemExOneAns}{AstraZeneca}
\newcommand{\TavilySemExOneNote}{Spacing only: "Astra Zeneca" == "AstraZeneca"}
\newcommand{\TavilySemExTwoGold}{US\$120,000}
\newcommand{\TavilySemExTwoAns}{\$120,000}
\newcommand{\TavilySemExTwoNote}{Same amount: "\$120,000" == "US\$120,000"}
\newcommand{\TavilySemExThreeGold}{3 players}
\newcommand{\TavilySemExThreeAns}{3}
\newcommand{\TavilySemExThreeNote}{Same value: "3" == gold "3 players"}
\newcommand{\FirecrawlSemExOneGold}{Astra Zeneca}
\newcommand{\FirecrawlSemExOneAns}{AstraZeneca}
\newcommand{\FirecrawlSemExOneNote}{Spacing only: "Astra Zeneca" == "AstraZeneca"}
\newcommand{\FirecrawlSemExTwoGold}{Bohemian Rhapsody}
\newcommand{\FirecrawlSemExTwoAns}{Bohemian Rhapsody, 9,948,386 viewers}
\newcommand{\FirecrawlSemExTwoNote}{Correct entity + extra detail: "Bohemian Rhapsody, 9,948,386 viewers"}
\newcommand{\FirecrawlSemExThreeGold}{16 years}
\newcommand{\FirecrawlSemExThreeAns}{16 years old}
\newcommand{\FirecrawlSemExThreeNote}{Same value: "16 years old" == "16 years"}


\newcommand{\smart}{\textsc{smart}}
\newcommand{\missed}{\textsc{missed}}
\newcommand{\blind}{\textsc{blind}}
\newcommand{\noop}{\textsc{no-op}}
\newcommand{\containsGold}{\texttt{contains\_gold\_answer}}
\newcommand{\contraGold}{\texttt{contradicts\_gold\_answer}}
\newcommand{\goldSnip}{\texttt{gold\_answer\_in\_snippets}}
\newcommand{\goldPage}{\texttt{gold\_answer\_in\_extracted\_page}}

\title{Search APIs Are Decision Surfaces: A Controlled Study of Tool-Using QA Agents}

\author{Anonymous ACL submission}

\begin{document}
\maketitle

\begin{abstract}
Search-augmented language-model agents are often built as if commercial search APIs were interchangeable retrieval backends. We test that assumption by freezing the agent, prompt, tools, and page-fetch backend, then swapping only the search provider. On \NQueries{} hard web-search questions, Brave, Tavily, and Firecrawl produce nearly indistinguishable exact-match accuracy (\BraveEM{}, \TavilyEM{}, and \FirecrawlEM{} out of 100). That headline hides large differences in the agent's internal retrieval economy. We introduce a per-URL oracle that labels every URL visible to the agent---not only URLs it fetched---and analyze search APIs as \emph{decision surfaces}: the ranked snippets and URLs that determine whether an agent answers, searches again, or spends fetch tokens. The same agent sees visible gold-supporting evidence on 33, 24, and 19 queries across the three providers; its four-way decision partition shifts from rich-snippet reading (Brave), to rank-one exploitation (Tavily), to broad blind exploration (Firecrawl). A surface-only contamination metric, the contradict-to-gold URL ratio, varies from 0.92 to 2.59 without looking at the model's answer. Provider choice is therefore not only a recall choice; it is a retrieval-budget allocation choice.
\end{abstract}

\section{Introduction}

A tool-using question-answering agent does not consume a search index directly. It consumes a \emph{surface}: titles, URLs, snippets, optional provider metadata, and ranks. Before the agent opens a page, those fields are the only evidence it has for deciding whether to answer, reformulate the query, or spend fetch tokens. This paper asks whether commercial search APIs can be treated as interchangeable surfaces for such agents.

The intuitive evaluation would compare final answer accuracy. We find that this view is misleading. In our controlled experiment, three production providers reach almost the same exact-match (EM) score: \BraveEM{}/100 for Brave, \TavilyEM{}/100 for Tavily, and \FirecrawlEM{}/100 for Firecrawl. But the trajectories that produce those scores are structurally different. One provider exposes more supporting evidence directly in snippets. Another concentrates support at rank one, making the agent's top-result fetches unusually productive. A third offers a sparse snippet surface but induces broad exploration and wins many queries outside the judged visible-support ceiling.

We call this phenomenon a \textbf{decision surface} effect. A search API is not merely a candidate generator for the answer model; it is the pre-fetch interface on which the model's policy is executed. Two providers can produce similar downstream EM while changing the conditional value of a fetch, the probability that a visible gold-bearing URL is opened, and the probability that contradicting snippets contaminate the run.

\paragraph{Contributions.} We make four contributions.
First, we define an agentic evaluation protocol that freezes the LLM agent and page fetcher while varying only the commercial search provider. Second, we introduce a per-URL oracle: an LLM judge labels every URL visible to the agent with support, contradiction, model-answer support, and garbage fields. Third, we define a four-cell decision partition---\smart{}, \missed{}, \blind{}, and \noop{}---that converts per-URL labels into provider-specific fetch-policy diagnostics. Fourth, we report an answer-independent surface-quality metric, the contradict-to-gold ratio, and show that it varies by 2.8$\times$ across providers even when EM does not.

\begin{figure*}[t]
\centering
\resizebox{0.98\textwidth}{!}{% Generated from Graphviz DOT by figures/make_figures.py; do not edit by hand.
\definecolor{slateink}{HTML}{334155}
\definecolor{slateedge}{HTML}{475569}
\definecolor{slatemid}{HTML}{64748B}
\definecolor{slatefill}{HTML}{F8FAFC}
\definecolor{orangefill}{HTML}{FFF7ED}
\definecolor{orangeedge}{HTML}{C2410C}
\definecolor{redfill}{HTML}{FEE2E2}
\definecolor{rededge}{HTML}{DC2626}
\definecolor{bluefill}{HTML}{DBEAFE}
\definecolor{blueedge}{HTML}{2563EB}
\definecolor{greenfill}{HTML}{DCFCE7}
\definecolor{greenedge}{HTML}{16A34A}
\definecolor{greenlight}{HTML}{F0FDF4}
\definecolor{greendark}{HTML}{15803D}
\definecolor{indigofill}{HTML}{EEF2FF}
\definecolor{indigoedge}{HTML}{4F46E5}
\definecolor{cyanfill}{HTML}{ECFEFF}
\definecolor{cyanedge}{HTML}{0891B2}
\definecolor{slatefillb}{HTML}{F1F5F9}
\definecolor{purplefill}{HTML}{FAE8FF}
\definecolor{purpleedge}{HTML}{A21CAF}
\definecolor{amberfill}{HTML}{FFFBEB}
\definecolor{amberedge}{HTML}{B45309}
\definecolor{pinkfill}{HTML}{FFF1F2}
\definecolor{bluelight}{HTML}{EFF6FF}
\definecolor{amberfillb}{HTML}{FEF3C7}
\definecolor{amberedgeb}{HTML}{D97706}
\definecolor{skyfill}{HTML}{E0F2FE}
\definecolor{skyedge}{HTML}{0284C7}
\definecolor{yellowfill}{HTML}{FEF9C3}
\definecolor{yellowedge}{HTML}{CA8A04}
\definecolor{orangefillb}{HTML}{FFEDD5}
\definecolor{orangeedgeb}{HTML}{EA580C}
\definecolor{grayfill}{HTML}{E2E8F0}
\begin{tikzpicture}[x=0.52in,y=0.52in,>=latex]
\tikzset{gvnode/.style={draw, rounded corners=3pt, align=center, font=\scriptsize, inner sep=2.8pt, line width=0.35pt}}
\tikzset{gvedge/.style={->, line width=0.35pt, color=slateedge}}
\draw[gvedge] (0.583,1.686) -- (0.583,1.997) -- (0.583,2.431) -- (0.583,2.431) -- (0.583,2.431) -- (1.950,2.431) -- (1.950,2.431);
\draw[gvedge] (1.171,1.375) -- (1.171,1.375) -- (1.947,1.375) -- (1.947,1.375);
\draw[gvedge] (0.583,1.064) -- (0.583,0.753) -- (0.583,0.319) -- (0.583,0.319) -- (0.583,0.319) -- (1.886,0.319) -- (1.886,0.319);
\draw[gvedge] (2.801,2.431) -- (3.473,2.431) -- (4.806,2.431) -- (4.806,2.431) -- (4.806,2.431) -- (4.806,1.783) -- (4.806,1.783);
\draw[gvedge] (2.800,1.375) -- (2.800,1.375) -- (3.583,1.375) -- (3.583,1.375);
\draw[gvedge] (2.866,0.361) -- (3.555,0.361) -- (4.806,0.361) -- (4.806,0.361) -- (4.806,0.361) -- (4.806,0.968) -- (4.806,0.968);
\draw[gvedge] (5.922,1.375) -- (5.922,1.375) -- (7.555,1.375) -- (7.555,1.375);
\node[font=\tiny, fill=white, inner sep=1pt, text=slatemid] at (7.555,1.375) {search surface};
\draw[gvedge] (9.823,1.607) -- (9.823,1.607) -- (11.733,1.607) -- (11.733,1.607);
\node[font=\tiny, fill=white, inner sep=1pt, text=slatemid] at (11.733,1.607) {selected URLs};
\draw[gvedge] (9.820,1.264) -- (10.684,1.264) -- (11.750,1.264) -- (11.750,1.264) -- (11.750,1.264) -- (11.750,0.865) -- (11.750,0.865);
\draw[gvedge] (11.831,1.643) -- (11.831,1.643) -- (9.920,1.643) -- (9.920,1.643);
\node[font=\tiny, fill=white, inner sep=1pt, text=slatemid] at (9.920,1.643) {markdown page};
\draw[gvedge] (13.181,0.766) -- (13.181,0.869) -- (13.181,0.958) -- (13.181,0.958) -- (13.181,0.958) -- (15.442,0.958) -- (15.442,0.958);
\node[font=\tiny, fill=white, inner sep=1pt, text=slatemid] at (13.181,0.958) {one row per visible URL};
\draw[gvedge] (13.266,0.361) -- (13.266,0.361) -- (21.623,0.361) -- (21.623,0.361);
\node[font=\tiny, fill=white, inner sep=1pt, text=slatemid] at (21.623,0.361) {actions + EM};
\draw[gvedge] (16.807,1.125) -- (16.807,1.125) -- (18.345,1.125) -- (18.345,1.125);
\draw[gvedge] (20.904,1.125) -- (21.766,1.125) -- (22.764,1.125) -- (22.764,1.125) -- (22.764,1.125) -- (22.764,0.667) -- (22.764,0.667);
\node[gvnode, minimum width=0.607in, minimum height=0.318in, fill=orangefill, draw=orangeedge, text=slateink] (q) at (0.583,1.375) {100\\SealQA-Hard\\queries};
\node[gvnode, minimum width=0.390in, minimum height=0.260in, fill=redfill, draw=rededge, text=slateink] (brave) at (2.424,2.431) {Brave};
\node[gvnode, minimum width=0.390in, minimum height=0.260in, fill=bluefill, draw=blueedge, text=slateink] (tavily) at (2.424,1.375) {Tavily};
\node[gvnode, minimum width=0.455in, minimum height=0.260in, fill=greenfill, draw=greenedge, text=slateink] (firecrawl) at (2.424,0.319) {Firecrawl};
\node[gvnode, minimum width=1.163in, minimum height=0.318in, fill=slatefill, draw=slateink, text=slateink] (surface) at (4.799,1.375) {provider\\decision surface\\(snippets + URLs + ranks)};
\node[gvnode, minimum width=1.127in, minimum height=0.318in, fill=indigofill, draw=indigoedge, text=slateink] (agent) at (8.736,1.375) {frozen GPT-5.4\\agent\\search\_web + fetch\_page};
\node[gvnode, minimum width=0.650in, minimum height=0.318in, fill=cyanfill, draw=cyanedge, text=slateink] (jina) at (12.458,1.875) {Jina Reader\\page fetcher\\fixed backend};
\node[gvnode, minimum width=0.838in, minimum height=0.318in, fill=slatefillb, draw=slateedge, text=slateink] (trace) at (12.458,0.458) {JSONL trace\\searches, fetches,\\final answer};
\node[gvnode, minimum width=0.657in, minimum height=0.318in, fill=purplefill, draw=purpleedge, text=slateink] (judge) at (16.174,1.153) {Kimi-K2.6\\per-URL judge\\temp. 0};
\node[gvnode, minimum width=1.278in, minimum height=0.260in, fill=greenlight, draw=greendark, text=slateink] (oracle) at (19.674,1.125) {visible-URL oracle\\gold / contradiction / garbage};
\node[gvnode, minimum width=1.076in, minimum height=0.260in, fill=amberfill, draw=amberedge, text=slateink] (metrics) at (22.757,0.319) {decision-surface metrics\\EM, buckets, rank, r\_c:g};
\end{tikzpicture}
}
\caption{Controlled experiment. The query set, agent, tools, prompt, iteration budget, and Jina Reader page-fetch backend are fixed. Only the search provider changes. The judge sees the same XML document surface seen by the agent before a fetch, plus fetched-page text when the agent opened that URL. Joining judge labels with trace actions yields decision-surface metrics. The diagram is generated from Graphviz DOT and rendered as TikZ for deterministic vector output.}
\label{fig:architecture}
\end{figure*}


\section{Related Work}
\label{sec:related}

\paragraph{Retrieval-augmented and tool-using QA.}
Retrieval-augmented generation evaluates how retrieved evidence improves knowledge-intensive generation \citep{lewis-etal-2020-retrieval,izacard-grave-2021-leveraging}. Tool-using agents extend this setting by letting the model choose when to call external tools \citep{yao2023react,schick-etal-2023-toolformer}. Our focus is not a new agent architecture; it is the effect of swapping the commercial search surface underneath a fixed agent.

\paragraph{Retrieval evaluation.}
Standard information-retrieval benchmarks emphasize static ranked-list metrics such as recall or nDCG over top-$k$ candidates \citep{thakur-etal-2021-beir}. Those metrics remain important, but they omit a key feature of agentic search: the model must decide from snippets whether to fetch, answer, or search again. Our per-URL oracle converts a ranked list into a decision surface by labeling the evidence state that was visible at each action point.

\paragraph{LLM judges and web-search benchmarks.}
LLM-as-judge methods are widely used to evaluate model outputs \citep{zheng2023judging}. We use a judge differently: it labels each retrieved document rather than only the final answer. The question set comes from \textsc{SealQA-Hard} \citep{sealqa2025}; the provider adapters use Brave, Tavily, Firecrawl, and Jina Reader APIs \citep{bravesearch-docs,tavily-docs,firecrawl-docs,jina-reader}.

\section{Task and Agent Setup}
\label{sec:setup}

\paragraph{Dataset.} We evaluate on a deterministic proportional stratified sample of \NQueries{} queries from the 254-row \textsc{SealQA-Hard} subset. The strata are \texttt{freshness} $\times$ \texttt{search\_results} $\times$ \texttt{topic}, using seed 20260509. The sample preserves the source distribution closely; Appendix~\ref{app:dataset} gives the allocation summary.

\paragraph{Agent.} The agent is a frozen GPT-5.4 tool user with a maximum of 10 iterations. Its two tools are \texttt{search\_web(query)}, which returns up to 10 ranked provider results, and \texttt{fetch\_page(document\_id)}, which opens a previously seen result. The prompt instructs the model to ground every final answer in retrieved evidence, to fetch when snippets are insufficient or conflicting, and to output only a short \texttt{FINAL ANSWER}. The prompt and tool schema are identical across provider conditions.

\paragraph{Providers.} We run the same 100 questions against Brave Web Search, Tavily Search, and Firecrawl Search. The provider adapters normalize each result into rank, title, URL, snippet, domain, and provider metadata. The Brave adapter stores provider-native \texttt{extra\_snippets}; Tavily and Firecrawl are configured with provider-side raw content or markdown disabled, so page text reaches the agent only through the shared \texttt{fetch\_page} tool.

\paragraph{Page fetcher.} All fetched pages are retrieved through Jina Reader and returned as markdown with fetch metadata. This isolates provider differences to the search surface: URL pool, ranking, snippet content, and metadata.

\section{Per-URL Oracle}
\label{sec:oracle}

For each provider-query trajectory, we judge every URL the agent saw. The judge is Kimi-K2.6 at temperature 0. It receives the question, gold answer, the model's final answer, and exactly one retrieved document rendered in the same XML-like format used by the answer model. For unfetched URLs, the document contains title, URL, domain, snippet, and provider metadata such as Brave \texttt{extra\_snippets}. For fetched URLs, the document additionally contains the extracted markdown page text.

The judge emits booleans for gold support, snippet-level support, page-level support, contradiction, model-answer support, and garbage, plus short evidence spans and confidence; Appendix~\ref{app:judge-schema} gives the exact JSON schema. We analyze \NJudgeValid{} valid judgments out of \NJudgeTotal{} total rows; \NJudgeInvalid{} invalid rows were timeouts or JSON/schema failures and are excluded. The valid rows split into \NJudgeSnippetValid{} snippet-only judgments and \NJudgePageValid{} page-visible judgments.

\paragraph{Observable ceilings.} We use two cautionary terms. A \emph{snippet-visible} support event means at least one judged snippet field contains enough evidence for the gold answer. A \emph{surface-visible support} event means at least one visible URL supports the gold answer either from snippets or from page text the agent actually fetched. The second is a lower bound on the true provider pool: for unfetched URLs, we do not know what the page would have contained.

\section{Metrics}
\label{sec:metrics}

Let a provider-query pair have a set of judged URLs $U$. For each URL $u$, the judge supplies support and contradiction labels, and the trace supplies whether the agent fetched $u$. We aggregate in three ways.

\paragraph{Headline answer quality.} We report exact match and token F1 against the SealQA gold answer. Because $n=100$ per provider, small EM differences are not treated as provider wins.

\paragraph{Decision partition.} Every provider-query pair belongs to exactly one bucket:
\begin{itemize}
    \item \smart{}: at least one visible URL supports the gold answer, and the agent fetched at least one such URL.
    \item \missed{}: at least one visible URL supports the gold answer, but the agent fetched none of the supporting URLs.
    \item \blind{}: no visible URL supports the gold answer, but the agent fetched at least one URL.
    \item \noop{}: no visible URL supports the gold answer, and the agent fetched nothing.
\end{itemize}
This partition is deliberately policy-aware: it separates provider surfaces that reveal evidence from agent policies that exploit or ignore that evidence.

\paragraph{Contradict-to-gold ratio.} On the snippet-only surface, define
\[
 r_{c:g}=\frac{\left|\{u \in U : \contraGold(u)\}\right|}{\left|\{u \in U : \containsGold(u)\}\right|}.
\]
The ratio does not use the model's final answer or fetched page text. It therefore measures surface contamination without circularly counting the model's mistakes.

\section{Results}
\label{sec:results}

\subsection{Final EM hides provider-specific policies}

\begin{table}[t]
\centering
\small
\begin{tabular}{lrrr}
\toprule
\textbf{Metric} & \textbf{Brave} & \textbf{Tavily} & \textbf{Firecrawl} \\
\midrule
Exact match / 100 & \BraveEM{} & \TavilyEM{} & \FirecrawlEM{} \\
Wilson 95\% CI & 14--30 & 14--30 & 16--32 \\
Token F1 & \BraveFone{} & \TavilyFone{} & \FirecrawlFone{} \\
Answered & 98 & 97 & 96 \\
\midrule
Avg. search calls & \BraveSearchAvg{} & \TavilySearchAvg{} & \FirecrawlSearchAvg{} \\
Avg. fetch calls & \BraveFetchAvg{} & \TavilyFetchAvg{} & \FirecrawlFetchAvg{} \\
Queries with $\geq$1 fetch & \BraveFetchedPct{} & \TavilyFetchedPct{} & \FirecrawlFetchedPct{} \\
Total agent tokens & \BraveTokensM{}M & \TavilyTokensM{}M & \FirecrawlTokensM{}M \\
\bottomrule
\end{tabular}
\caption{Headline trajectory metrics. End-to-end answer accuracy is nearly unchanged across providers, while search and fetch behavior differs.}
\label{tab:headline}
\end{table}

Table~\ref{tab:headline} is the result a standard end-to-end evaluation would emphasize: EM is 21--23 out of 100, with heavily overlapping Wilson intervals. A practitioner looking only at this table could conclude that provider choice is not important for answer accuracy. The remaining results show why that conclusion is incomplete.

\begin{figure*}[t]
\centering
\resizebox{0.98\textwidth}{!}{% Generated from Graphviz DOT by figures/make_figures.py; do not edit by hand.
\definecolor{slateink}{HTML}{334155}
\definecolor{slateedge}{HTML}{475569}
\definecolor{slatemid}{HTML}{64748B}
\definecolor{slatefill}{HTML}{F8FAFC}
\definecolor{orangefill}{HTML}{FFF7ED}
\definecolor{orangeedge}{HTML}{C2410C}
\definecolor{redfill}{HTML}{FEE2E2}
\definecolor{rededge}{HTML}{DC2626}
\definecolor{bluefill}{HTML}{DBEAFE}
\definecolor{blueedge}{HTML}{2563EB}
\definecolor{greenfill}{HTML}{DCFCE7}
\definecolor{greenedge}{HTML}{16A34A}
\definecolor{greenlight}{HTML}{F0FDF4}
\definecolor{greendark}{HTML}{15803D}
\definecolor{indigofill}{HTML}{EEF2FF}
\definecolor{indigoedge}{HTML}{4F46E5}
\definecolor{cyanfill}{HTML}{ECFEFF}
\definecolor{cyanedge}{HTML}{0891B2}
\definecolor{slatefillb}{HTML}{F1F5F9}
\definecolor{purplefill}{HTML}{FAE8FF}
\definecolor{purpleedge}{HTML}{A21CAF}
\definecolor{amberfill}{HTML}{FFFBEB}
\definecolor{amberedge}{HTML}{B45309}
\definecolor{pinkfill}{HTML}{FFF1F2}
\definecolor{bluelight}{HTML}{EFF6FF}
\definecolor{amberfillb}{HTML}{FEF3C7}
\definecolor{amberedgeb}{HTML}{D97706}
\definecolor{skyfill}{HTML}{E0F2FE}
\definecolor{skyedge}{HTML}{0284C7}
\definecolor{yellowfill}{HTML}{FEF9C3}
\definecolor{yellowedge}{HTML}{CA8A04}
\definecolor{orangefillb}{HTML}{FFEDD5}
\definecolor{orangeedgeb}{HTML}{EA580C}
\definecolor{grayfill}{HTML}{E2E8F0}
\begin{tikzpicture}[x=0.52in,y=0.52in,>=latex]
\tikzset{gvnode/.style={draw, rounded corners=3pt, align=center, font=\scriptsize, inner sep=2.8pt, line width=0.35pt}}
\tikzset{gvedge/.style={->, line width=0.35pt, color=slateedge}}
\draw[gvedge] (0.892,0.285) -- (0.892,0.285) -- (1.366,0.285) -- (1.366,0.285);
\draw[gvedge] (2.627,0.285) -- (2.627,0.285) -- (3.199,0.285) -- (3.199,0.285);
\draw[gvedge] (4.668,0.285) -- (4.668,0.285) -- (5.268,0.285) -- (5.268,0.285);
\draw[gvedge] (0.892,1.271) -- (0.892,1.271) -- (1.366,1.271) -- (1.366,1.271);
\draw[gvedge] (2.627,1.271) -- (2.627,1.271) -- (3.104,1.271) -- (3.104,1.271);
\draw[gvedge] (4.768,1.271) -- (4.768,1.271) -- (5.268,1.271) -- (5.268,1.271);
\draw[gvedge] (0.892,2.215) -- (0.892,2.215) -- (1.366,2.215) -- (1.366,2.215);
\draw[gvedge] (2.627,2.215) -- (2.627,2.215) -- (3.206,2.215) -- (3.206,2.215);
\draw[gvedge] (4.662,2.215) -- (4.662,2.215) -- (5.243,2.215) -- (5.243,2.215);
\node[gvnode, minimum width=0.462in, minimum height=0.260in, fill=redfill, draw=rededge, text=slateink] (b0) at (0.444,0.285) {Brave\\EM 21/100};
\node[gvnode, minimum width=0.614in, minimum height=0.260in, fill=pinkfill, draw=rededge, text=slateink] (b1) at (2.035,0.285) {visible support\\33 queries};
\node[gvnode, minimum width=0.722in, minimum height=0.296in, fill=pinkfill, draw=rededge, text=slateink] (b2) at (3.972,0.285) {large snippet-side\\support pool\\97 gold URLs};
\node[gvnode, minimum width=0.744in, minimum height=0.260in, fill=pinkfill, draw=rededge, text=slateink] (b3) at (6.062,0.285) {low contamination\\r\_c:g = 0.92};
\node[gvnode, minimum width=0.462in, minimum height=0.260in, fill=bluefill, draw=blueedge, text=slateink] (t0) at (0.444,1.271) {Tavily\\EM 21/100};
\node[gvnode, minimum width=0.614in, minimum height=0.260in, fill=bluelight, draw=blueedge, text=slateink] (t1) at (2.035,1.271) {visible support\\24 queries};
\node[gvnode, minimum width=0.823in, minimum height=0.260in, fill=bluelight, draw=blueedge, text=slateink] (t2) at (3.972,1.271) {rank-1 concentration\\15/31 gold URLs};
\node[gvnode, minimum width=0.744in, minimum height=0.260in, fill=bluelight, draw=blueedge, text=slateink] (t3) at (6.062,1.271) {fetching gold pays\\SMART EM 64\%};
\node[gvnode, minimum width=0.462in, minimum height=0.260in, fill=greenfill, draw=greenedge, text=slateink] (f0) at (0.444,2.215) {Firecrawl\\EM 23/100};
\node[gvnode, minimum width=0.614in, minimum height=0.260in, fill=greenlight, draw=greenedge, text=slateink] (f1) at (2.035,2.215) {visible support\\19 queries};
\node[gvnode, minimum width=0.715in, minimum height=0.260in, fill=greenlight, draw=greenedge, text=slateink] (f2) at (3.972,2.215) {broad exploration\\67 BLIND queries};
\node[gvnode, minimum width=0.773in, minimum height=0.260in, fill=greenlight, draw=greenedge, text=slateink] (f3) at (6.062,2.215) {high contamination\\r\_c:g = 2.59};
\end{tikzpicture}
}
\caption{Provider profiles under the same agent. Brave exposes the broadest visible-support set and the lowest contradict-to-gold ratio; Tavily concentrates support at rank one and makes fetching gold-supporting URLs high-payoff; Firecrawl reaches the highest EM despite the smallest visible-support set by inducing the largest blind-exploration bucket. Generated from Graphviz DOT.}
\label{fig:provider-profiles}
\end{figure*}

\subsection{The same EM comes from three different surfaces}

\begin{table}[t]
\centering
\small
\begin{tabular}{lrrr}
\toprule
\textbf{Surface / oracle statistic} & \textbf{Brave} & \textbf{Tavily} & \textbf{Firecrawl} \\
\midrule
Valid snippet-only URL rows & \BraveSnippetRows{} & \TavilySnippetRows{} & \FirecrawlSnippetRows{} \\
Valid page-visible URL rows & \BravePageRows{} & \TavilyPageRows{} & \FirecrawlPageRows{} \\
Queries with visible support & \BraveSupportVisibleQ{} & \TavilySupportVisibleQ{} & \FirecrawlSupportVisibleQ{} \\
Queries without visible support & \BraveNoSupportQ{} & \TavilyNoSupportQ{} & \FirecrawlNoSupportQ{} \\
Gold-supporting snippet URLs & \BraveGoldURLs{} & \TavilyGoldURLs{} & \FirecrawlGoldURLs{} \\
Contradicting snippet URLs & \BraveContraURLs{} & \TavilyContraURLs{} & \FirecrawlContraURLs{} \\
Gold URLs at rank 1 & \BraveRankOnePct{} & \textbf{\TavilyRankOnePct{}} & \FirecrawlRankOnePct{} \\
\bottomrule
\end{tabular}
\caption{Oracle-visible support and contamination. ``Visible support'' is a lower bound over judged snippets plus fetched pages; unfetched pages are not page-judged.}
\label{tab:surface}
\end{table}

The provider surfaces differ before the agent answers. Brave has the largest support-visible ceiling (33 queries) and the largest count of snippet-level gold-supporting URLs. Tavily has fewer gold-supporting snippet URLs, but \TavilyRankOneGold{} of them are at rank one, so the agent's default tendency to open high-ranked results is unusually well aligned with Tavily's surface. Firecrawl has the smallest visible-support ceiling but the largest unsupported region; its final EM depends more on exploration than on the judge-visible support surface.

\subsection{The fetch policy changes sign across providers}

\begin{figure}[t]
\centering
\resizebox{\columnwidth}{!}{% Generated from Graphviz DOT by figures/make_figures.py; do not edit by hand.
\definecolor{slateink}{HTML}{334155}
\definecolor{slateedge}{HTML}{475569}
\definecolor{slatemid}{HTML}{64748B}
\definecolor{slatefill}{HTML}{F8FAFC}
\definecolor{orangefill}{HTML}{FFF7ED}
\definecolor{orangeedge}{HTML}{C2410C}
\definecolor{redfill}{HTML}{FEE2E2}
\definecolor{rededge}{HTML}{DC2626}
\definecolor{bluefill}{HTML}{DBEAFE}
\definecolor{blueedge}{HTML}{2563EB}
\definecolor{greenfill}{HTML}{DCFCE7}
\definecolor{greenedge}{HTML}{16A34A}
\definecolor{greenlight}{HTML}{F0FDF4}
\definecolor{greendark}{HTML}{15803D}
\definecolor{indigofill}{HTML}{EEF2FF}
\definecolor{indigoedge}{HTML}{4F46E5}
\definecolor{cyanfill}{HTML}{ECFEFF}
\definecolor{cyanedge}{HTML}{0891B2}
\definecolor{slatefillb}{HTML}{F1F5F9}
\definecolor{purplefill}{HTML}{FAE8FF}
\definecolor{purpleedge}{HTML}{A21CAF}
\definecolor{amberfill}{HTML}{FFFBEB}
\definecolor{amberedge}{HTML}{B45309}
\definecolor{pinkfill}{HTML}{FFF1F2}
\definecolor{bluelight}{HTML}{EFF6FF}
\definecolor{amberfillb}{HTML}{FEF3C7}
\definecolor{amberedgeb}{HTML}{D97706}
\definecolor{skyfill}{HTML}{E0F2FE}
\definecolor{skyedge}{HTML}{0284C7}
\definecolor{yellowfill}{HTML}{FEF9C3}
\definecolor{yellowedge}{HTML}{CA8A04}
\definecolor{orangefillb}{HTML}{FFEDD5}
\definecolor{orangeedgeb}{HTML}{EA580C}
\definecolor{grayfill}{HTML}{E2E8F0}
\begin{tikzpicture}[x=0.52in,y=0.52in,>=latex]
\tikzset{gvnode/.style={draw, rounded corners=3pt, align=center, font=\scriptsize, inner sep=2.8pt, line width=0.35pt}}
\tikzset{gvedge/.style={->, line width=0.35pt, color=slateedge}}
\draw[gvedge] (5.410,4.167) -- (5.410,4.167) -- (5.410,3.677) -- (5.410,3.677);
\draw[gvedge] (4.625,3.024) -- (4.625,3.024) -- (4.625,1.915) -- (4.625,1.915);
\node[font=\tiny, fill=white, inner sep=1pt, text=slatemid] at (4.625,1.915) {yes};
\draw[gvedge] (6.201,3.030) -- (6.201,3.030) -- (6.201,1.911) -- (6.201,1.911);
\node[font=\tiny, fill=white, inner sep=1pt, text=slatemid] at (6.201,1.911) {no};
\draw[gvedge] (2.826,1.667) -- (2.030,1.667) -- (1.104,1.667) -- (1.104,1.667) -- (1.104,1.667) -- (1.104,0.647) -- (1.104,0.647);
\node[font=\tiny, fill=white, inner sep=1pt, text=slatemid] at (1.104,1.667) {yes};
\draw[gvedge] (3.951,1.248) -- (3.951,1.248) -- (3.951,0.647) -- (3.951,0.647);
\node[font=\tiny, fill=white, inner sep=1pt, text=slatemid] at (3.951,0.647) {no};
\draw[gvedge] (6.882,1.248) -- (6.882,1.248) -- (6.882,0.647) -- (6.882,0.647);
\node[font=\tiny, fill=white, inner sep=1pt, text=slatemid] at (6.882,0.647) {yes};
\draw[gvedge] (8.011,1.667) -- (8.836,1.667) -- (9.812,1.667) -- (9.812,1.667) -- (9.812,1.667) -- (9.812,0.647) -- (9.812,0.647);
\node[font=\tiny, fill=white, inner sep=1pt, text=slatemid] at (9.812,1.667) {no};
\node[gvnode, minimum width=1.343in, minimum height=0.260in, fill=slatefill, draw=slateink, text=slateink] (start) at (5.410,4.417) {For each provider-query pair\\join judge labels with agent actions};
\node[gvnode, diamond, aspect=2.1, fill=amberfillb, draw=amberedgeb, text=slateink] (gold) at (5.410,3.181) {Any visible URL\\contains gold?};
\node[gvnode, diamond, aspect=2.1, fill=greenfill, draw=greenedge, text=slateink] (fetched_gold) at (3.951,1.667) {Agent fetched\\a gold URL?};
\node[gvnode, diamond, aspect=2.1, fill=skyfill, draw=skyedge, text=slateink] (fetched_any) at (6.882,1.667) {Agent fetched\\any URL?};
\node[gvnode, minimum width=1.148in, minimum height=0.296in, fill=greenfill, draw=greenedge, text=slateink] (smart) at (1.104,0.285) {SMART\\visible support + fetched it\\Brave 8, Tavily 11, Firecrawl 7};
\node[gvnode, minimum width=1.235in, minimum height=0.296in, fill=yellowfill, draw=yellowedge, text=slateink] (missed) at (3.951,0.285) {MISSED\\visible support + did not fetch it\\Brave 25, Tavily 13, Firecrawl 12};
\node[gvnode, minimum width=1.235in, minimum height=0.296in, fill=orangefillb, draw=orangeedgeb, text=slateink] (blind) at (6.882,0.285) {BLIND\\no visible support + fetched\\Brave 51, Tavily 55, Firecrawl 67};
\node[gvnode, minimum width=1.235in, minimum height=0.296in, fill=grayfill, draw=slateedge, text=slateink] (noop) at (9.812,0.285) {NO-OP\\no visible support + no fetch\\Brave 16, Tavily 21, Firecrawl 14};
\end{tikzpicture}
}
\caption{The four-cell decision partition. A provider changes both the probability that support is visible and the probability that the agent fetches a supporting URL when one is visible. Generated from Graphviz DOT.}
\label{fig:partition}
\end{figure}

\begin{table}[t]
\centering
\small
\begin{tabular}{lccc}
\toprule
\textbf{Bucket} & \textbf{Brave} & \textbf{Tavily} & \textbf{Firecrawl} \\
\midrule
\smart{} & \BraveSmartN{} / 38\% & \TavilySmartN{} / \textbf{64\%} & \FirecrawlSmartN{} / 29\% \\
\missed{} & \BraveMissedN{} / 36\% & \TavilyMissedN{} / 38\% & \FirecrawlMissedN{} / \textbf{42\%} \\
\blind{} & \BraveBlindN{} / 18\% & \TavilyBlindN{} / 16\% & \FirecrawlBlindN{} / \textbf{24\%} \\
\noop{} & \BraveNoopN{} / 0\% & \TavilyNoopN{} / 0\% & \FirecrawlNoopN{} / 0\% \\
\midrule
\smart{}--\missed{} gap & +2 pp & \textbf{+26 pp} & -13 pp \\
\bottomrule
\end{tabular}
\caption{Bucket size and EM rate. The same fetch interface has positive, near-zero, and negative observed return depending on the provider surface. The Firecrawl gap has small $n$ and should be read as exploratory.}
\label{tab:buckets}
\end{table}

Table~\ref{tab:buckets} shows that fetch decisions are not provider-neutral. On Tavily, opening a gold-supporting URL is associated with a large EM lift: \smart{} queries are correct 64\% of the time, compared with 38\% for \missed{} queries. On Brave, the difference is negligible because many answers are already readable from snippets or extra snippets. On Firecrawl, the \smart{} cell is small and lower-EM than \missed{}; Firecrawl's correct answers concentrate in the \blind{} bucket, suggesting that the agent benefits from broad exploration of a sparse surface.

The key point is not that one static policy wins. It is that the same policy has a different return curve on each provider. Treating providers as interchangeable therefore hides a budget-allocation decision.

\subsection{Surface contamination is measurable without model answers}

\begin{table}[t]
\centering
\small
\begin{tabular}{lrrr}
\toprule
\textbf{Provider} & \textbf{Contradict} & \textbf{Gold} & $\boldsymbol{r_{c:g}}$ \\
\midrule
Brave & \BraveContraURLs{} & \BraveGoldURLs{} & \textbf{\BraveContraRatio{}} \\
Tavily & \TavilyContraURLs{} & \TavilyGoldURLs{} & \TavilyContraRatio{} \\
Firecrawl & \FirecrawlContraURLs{} & \FirecrawlGoldURLs{} & \textbf{\FirecrawlContraRatio{}} \\
\bottomrule
\end{tabular}
\caption{Contradict-to-gold ratio on snippet-only URL rows. The metric is answer-independent: it does not condition on whether the agent was correct.}
\label{tab:contamination}
\end{table}

The contradict-to-gold ratio ranges from 0.92 for Brave to 2.59 for Firecrawl. This is the largest provider spread we observe on a surface-only metric. It also avoids a common circularity: counting documents that support the model's wrong answer can make contamination appear high simply because the model made a mistake. The \contraGold{} label depends only on the question, gold answer, and document surface.

\subsection{Answer text availability is not enough}

The deterministic trace metrics also show why retrieval evaluation cannot stop at answer-string availability. The gold answer string appears somewhere in retrieved snippets, extra snippets, or fetched pages on 78 Brave queries, 75 Tavily queries, and 76 Firecrawl queries, yet EM remains near 21--23. Thus the dominant failure mode is not simply failure to retrieve an answer string. It includes deciding which surface to trust, whether to fetch, how to resolve contradictions, and how to convert page evidence into the exact requested granularity.

\section{Discussion}
\label{sec:discussion}

\paragraph{Provider choice is policy choice.} A provider surface changes the prior over which action is profitable. Brave rewards reading a richer snippet surface; Tavily rewards fetching high-ranked results; Firecrawl rewards wider exploration. A deployed agent should not necessarily use the same fetch threshold or rank cutoff for all providers.

\paragraph{Top-$k$ relevance misses the mechanism.} A top-$k$ URL metric can tell us whether relevant pages are somewhere in the pool. It cannot tell us whether the snippet is sufficient, whether the page needs opening, whether contradicting snippets surround the answer, or whether the agent's fetch heuristic will select the useful URL. Decision-surface metrics complement top-$k$ relevance by making these pre-fetch choices observable.

\paragraph{Design implication.} Search APIs for agents should optimize not only traditional relevance but also actionability: calibrated snippets, fewer unsupported contradictions, stable document identifiers, useful rank concentration, and metadata that tells the agent when opening the page is likely to pay off.

\section{Threats to Validity}
\label{sec:limitations}

\paragraph{Single agent and single judge.} The decision partition is induced by one GPT-5.4 agent and one Kimi-K2.6 judge. Different agents may fetch more aggressively, and a second judge could shift boundary cases.

\paragraph{Observable-pool lower bound.} We page-judge only URLs the agent fetched. An unfetched URL may contain gold on the page even when its snippet does not. Therefore visible support is a lower bound on true provider pool support, not a full recall estimate.

\paragraph{Snippet asymmetry is real but confounded.} Brave exposes \texttt{extra\_snippets}; Tavily and Firecrawl do not under our configuration. This is a provider feature visible to the agent, so it belongs in the surface, but it complicates any claim that compares snippet text length rather than policy outcome.

\paragraph{Observational policy analysis.} We observe the agent's actions under each provider. We do not replay the same agent hidden state with swapped result lists, nor do we intervene on ranks or snippets. The current study identifies decision-surface effects; a causal replay study is the natural next step.

\paragraph{Scale.} With 100 queries, EM intervals are wide. We therefore emphasize structural diagnostics and report cell sizes, not provider leaderboards.

\section{Conclusion}

Commercial search APIs are not interchangeable black boxes for tool-using agents. In our controlled study, final answer accuracy barely moves, but the internal retrieval economy changes sharply. Search APIs should be evaluated as decision surfaces: ranked evidence interfaces that shape when an agent answers, fetches, explores blindly, or becomes contaminated by contradictions. The practical question is not only ``which provider retrieves relevant URLs?'' but ``which provider induces the right retrieval-budget policy for this agent?''

\section*{Limitations}
The main limitations are the single-agent/single-judge design, the lower-bound nature of page-visible support, the provider-specific snippet asymmetry, the observational rather than causal policy analysis, and the 100-query scale. These limitations constrain effect-size claims but do not remove the central observation that similar EM can arise from different decision policies.

\section*{Ethics Statement}
This work evaluates commercial search APIs and LLM agents on public web-search QA data. We do not rank providers as generally better or worse; we report provider-specific behavior under a fixed research protocol. The study may help practitioners reduce unnecessary page fetches and better detect contradictory evidence, but the same measurements could be used to overfit agents to a benchmark. We therefore release metric definitions and emphasize replication on additional datasets and agents before deployment decisions.

\section*{Reproducibility Statement}
All numbers are deterministic functions of the committed query sample, trace JSONL files, judge JSONL files, and provider-comparison summaries. The raw trace and judge JSONLs are stored through Git LFS; run \texttt{git lfs pull} before executing \texttt{paper/figures/make\_figures.py}. The script regenerates the Graphviz-derived TikZ figures, the number macros, and an audit file without making model or provider API calls.

\bibliography{refs}
\bibliographystyle{acl_natbib}

\appendix

\section{Dataset Sample}
\label{app:dataset}
The sample file records \texttt{dataset=vtllms/sealqa}, \texttt{subset=seal\_hard}, sample size 100, source size 254, seed 20260509, and strata fields \texttt{freshness}, \texttt{search\_results}, and \texttt{topic}. The sampling diagnostics are retained with per-cell source and sample counts.

\section{Provider Requests}
\label{app:provider-requests}
Brave uses \texttt{GET https://api.search.brave.com/res/v1/web/search} with \texttt{count=10}, \texttt{country=US}, English language settings, \texttt{safesearch=moderate}, \texttt{spellcheck=true}, \texttt{text\_decorations=false}, and \texttt{result\_filter=web}. Tavily uses \texttt{POST https://api.tavily.com/search} with \texttt{search\_depth=basic}, \texttt{topic=general}, \texttt{max\_results=10}, and provider-side answer/image/raw-content options disabled. Firecrawl uses \texttt{POST https://api.firecrawl.dev/v2/search} with \texttt{limit=10}, \texttt{sources=[web]}, \texttt{country=US}, \texttt{timeout=60000}, and no markdown scraping.

\section{Wilson Intervals for Decision Buckets}
\label{app:wilson}
\begin{table}[h]
\centering
\small
\begin{tabular}{lccc}
\toprule
\textbf{Bucket} & \textbf{Brave} & \textbf{Tavily} & \textbf{Firecrawl} \\
\midrule
\smart{} & 3/8: 38\% [14--69] & 7/11: 64\% [35--85] & 2/7: 29\% [8--64] \\
\missed{} & 9/25: 36\% [20--55] & 5/13: 38\% [18--64] & 5/12: 42\% [19--68] \\
\blind{} & 9/51: 18\% [10--30] & 9/55: 16\% [9--28] & 16/67: 24\% [15--35] \\
\noop{} & 0/16: 0\% [0--19] & 0/21: 0\% [0--15] & 0/14: 0\% [0--22] \\
\bottomrule
\end{tabular}
\caption{Wilson 95\% intervals for EM rates in the decision partition.}
\end{table}

\section{Judge Schema}
\label{app:judge-schema}
\begin{small}
\begin{verbatim}
{
  "contains_gold_answer": boolean,
  "gold_answer_in_snippets": boolean,
  "gold_answer_in_extracted_page": boolean,
  "supports_model_answer": boolean,
  "contradicts_gold_answer": boolean,
  "is_garbage": boolean,
  "garbage_reason": string,
  "answer_span": string,
  "gold_snippet_span": string,
  "gold_extracted_page_span": string,
  "confidence": number
}
\end{verbatim}
\end{small}

\end{document}
